\documentclass[french]{book}

\usepackage[utf8]{inputenc}
\usepackage[T1]{fontenc}
\usepackage{lmodern}
\usepackage{geometry}
\usepackage{graphicx}
\usepackage{babel}
\usepackage{amsmath}
\usepackage{amsthm}
\usepackage{amssymb}
\usepackage{hyperref}
\usepackage{stmaryrd}
\usepackage{enumitem}
\usepackage[most]{tcolorbox}
\usepackage{dsfont}
\usepackage{multirow}
\usepackage{etoc}
\usepackage{titlesec}
\usepackage{esint}
\usepackage{cancel}
\usepackage{mathdots}
\usepackage{indentfirst}

\setlength{\parindent}{0pt}

\setlist[itemize]{label=$\bullet$}
\setlist{before=\leavevmode}

\renewcommand{\P}{\mathbb{P}}
\newcommand{\N}{\mathbb{N}}
\newcommand{\Z}{\mathbb{Z}}
\newcommand{\Q}{\mathbb{Q}}
\newcommand{\R}{\mathbb{R}}
\newcommand{\C}{\mathbb{C}}
\newcommand{\K}{\mathbb{K}}
\renewcommand{\L}{\mathbb{L}}
\newcommand{\U}{\mathbb{U}}
\newcommand{\st}{^*}
\newcommand{\tq}{\middle|}
\newcommand{\inv}{^{-1}}
\newcommand{\E}{\mathbb{E}}
\newcommand{\F}{\mathbb{F}}
\newcommand{\V}{\mathbb{V}}
\renewcommand{\ker}{\text{Ker}}
\newcommand{\im}{\text{Im}}
\newcommand{\card}{\text{Card}}
\newcommand{\Tr}{\text{Tr}}
\newcommand{\Vect}{\text{Vect}}
\newcommand{\rg}{\text{rg}}
\newcommand{\vertiii}[1]{{\left\vert\kern-0.25ex\left\vert\kern-0.25ex\left\vert #1 \right\vert\kern-0.25ex\right\vert\kern-0.25ex\right\vert}}
\renewcommand{\o}{\text{o}}
\renewcommand{\O}{\text{O}}
\renewcommand{\d}{\text{d}}
\newcommand{\Id}{\text{Id}}
\newcommand{\D}{\text{D}}

\theoremstyle{definition}
\newtheorem{ex}{Exercice}[chapter]
\newtheorem{pro}{Problème}[chapter]

\theoremstyle{remark}
\newtheorem*{sol}{Solution}
\newtheorem*{rmq}{Remarque}

\title{Exercices intéressants}

\begin{document}

\chapter{Polynômes et fractions rationnelles}

\section{Problèmes}

\begin{pro}
	Pour une fraction rationnelle non nulle $F\in\C(X)$, on note $n(F)$ le nombre de ses racines (distinctes), et $p(F)$ le nombre de ses pôles.
	\begin{enumerate}
		\item Soit $F\in\C(X)$ non constante. Montrer que $n(F)+n(F-1)+p(F)>a$, où $a$ désigne le degré du numérateur de $F$ dans son écriture sous forme irréductible. On pourra commencer par le cas où $F$ est un polynôme.
		\item Soit $P$ et $Q$ deux polynômes non constants dans $\C[X]$ tels que $P\inv(\{0\})=Q\inv(\{0\})$ et $P\inv(\{1\})=Q\inv(\{1\})$. Montrer que $P=Q$.
		\item Théorème de Mason : soit $A$, $B$ et $C$ trois polynômes non nuls premiers entre eux dans leur ensemble, non tous constants et tels que $A+B+C=0$. On note $n$ le maximum des degrés de $A$, $B$ et $C$. Montrer que $ABC$ possède au moins $n+1$ racines distinctes.
		\item Soit $m\geqslant 3$ entier et $(A,B,C)\in\C[X]^3$ un triplet de polynômes premiers entre eux dans leur ensemble tel que $A^m+B^m=C^m$. Montrer que $A$, $B$ et $C$ sont tous constants.
	\end{enumerate}
\end{pro}

\section{Exercices}

\begin{ex}[Critère d'Eisenstein]
	Soit $p$ un nombre premier et $P\in\Z[X]$ qu'on écrit sous la forme $$P=\sum_{k=0}^{n}a_kX^k$$ Supposons que l'on ait :
	\begin{itemize}
		\item $\forall k\in\llbracket0,n-1\rrbracket,\ p\mid a_k$ ;
		\item $p^2\nmid a_0$.
		\item $p\nmid a_n$
	\end{itemize}
	Montrer que $P$ est irréductible dans $\Z[X]$.
\end{ex}

\begin{ex}[Polynômes à valeurs dans les nombres premiers]
	On note $\mathcal{P}$ l'ensemble des nombres premiers. Trouver tous les polynômes $A \in \Z[X]$ tels que $\forall n \in \N, \ A(n) \in \mathcal{P}$.
\end{ex}

\chapter{Révisions sur les espaces vectoriels}

\section{Problèmes}

\section{Exercices}

\begin{ex}
	Soit $n\in\N\st$. On se donne $X_0,\ldots,X_n$ des parties non vides de $\llbracket1,n\rrbracket$. Montrer qu'il existe $I$ et $J$ deux parties disjointes non vides de $\llbracket0,n\rrbracket$ telles que : $$\bigcup_{i\in I} X_i=\bigcup_{i\in J} X_j$$
\end{ex}

\begin{ex}[Familles positivement génératrices]
	Soit $E$ un $\R$-espace vectoriel et $(e_i)_{i\in I}$ une famille de vecteurs de $E$. On dit que $(e_i)_{i\in I}$ est positivement génératrice lorsque tout vecteur de $E$ est combinaison linéaire à coefficients strictement positifs des $e_i$. Dans tout l'exercice, on suppose $E$ de dimension finie $n\in\N\st$ et on se donne $(e_1,\ldots,e_p)\in E^p$ une famille positivement génératrice.
	\begin{enumerate}
		\item Montrer : $p\geqslant n+1$.
		\item Donner un exemple où on peut avoir $p=n+1$.
		\item On suppose de plus $p\geqslant 2n+1$. Montrer que $(e_1,\ldots,e_p)$ possède une sous-famille stricte qui est toujours positivement génératrice.
		\item Donner un exemple où $p=2n$ et où aucune sous-famille stricte n'est positivement génératrice.
	\end{enumerate}
\end{ex}
\begin{sol}
	\begin{enumerate}
		\item Si une famille est positivement génératrice, alors elle est génératrice donc on a déjà $p\geqslant n$. Par l'absurde, si $p=n$, alors $(e_1,\ldots,e_p)$ est une base de $E$. Mais alors, le vecteur $-(e_1+\cdots+e_p)$ ne peut s'écrire comme combinaison linéaire à coefficients strictement positifs des $e_i$, ce qui est absurde. Donc $p\geqslant n+1$.
		\item Notons $(f_1,\ldots,f_n)$ une base de $E$. Prendre :
		$$(e_1,\ldots,e_n,e_{n+1})=(f_1,\ldots,f_n,-(f_1+\cdots+f_n))$$ convient.
		\item En utilisant le même genre d'idées que dans la question précédente, exprimer un vecteur comme combinaison linéaire à coefficients strictement positifs des autres. Pour cela, extraire une base de ces $2n+1$ vecteurs, puis remarquer que les $n+1$ derniers forment une famille liée.
		\item En gardant les notations de la question $2$, la famille suivante devrait fonctionner :
		$$(e_1,\ldots,e_n,e_{n+1},\ldots,e_{2n})=(f_1,\ldots,f_n,-f_1,\ldots,-f_n).$$
	\end{enumerate}
\end{sol}

\begin{ex}[Théorème de Maschke]
	Soit $E$ un espace vectoriel de dimension finie, $G$ un sous-groupe fini de $\mathcal{GL}(E)$ et $F$ un sous-espace vectoriel de $E$ stable par tous les éléments de $G$. Montrer que $F$ admet un supplémentaire stable par tous les éléments de $G$. On pourra considérer la moyenne des conjugués d'un projecteur sur $F$ par les éléments de $G$.
\end{ex}
\begin{sol}
	Notons $F'$ un supplémentaire de $F$ dans $E$ et $p$ le projecteur sur $F$ parallèlement à $F'$. On pose :
	$$q=\frac{1}{\card(G)}\sum_{g\in G}g\inv\circ p\circ g.$$
	Montrons que $q$ est un projecteur d'image $F$.\\
	Déjà, $q\circ p=p$ donc $F=\text{Im}(p)\subset\text{Im}(q)$. \\
	Ensuite, si $h\in G$, par changement d'indice, on a $h\circ q=q\circ h$. Ainsi :
	\begin{align*}
		q^2&=\frac{1}{\card(G)}\sum_{h\in G}q\circ h\inv \circ p\circ h\\
		&=\frac{1}{\card(G)}\sum_{h\in G}h\inv\circ q\circ p\circ  h\\
		&=\frac{1}{\card(G)}\sum_{h\in G}h\inv\circ p\circ  h\\
		&=q.
	\end{align*}
	$q$ a même rang (car même trace en dimension finie) que $p$, et son image contient $F$, donc $q$ est un projecteur d'image $F$. Notons $F''=\ker(q)$. Par changement d'indice, $F''$ est stable par tous les éléments de $G$, et c'est un supplémentaire de $F$ dans $E$, donc il convient.
\end{sol}

\chapter{Révisions sur les matrices}

\section{Problèmes}

\section{Exercices}

\begin{ex}
	Soit $A\in\mathcal{M}_n(\K)$ non scalaire. Montrer que pour tout $(a_1,\ldots,a_n)\in\K^n$ tel que $\Tr(A)=a_1+\cdots+a_n$, il existe une matrice semblable à $A$ dont les coefficients diagonaux sont $a_1,\ldots,a_n$.
\end{ex}

\begin{ex}
	Soit $n\in\N\st$. Déterminer le $k\in\N$ maximal tel qu'il existe des parties $E_1,\ldots,E_k$ de $\llbracket 1,n\rrbracket$ telles que :
	\begin{itemize}
		\item $\card(E_i)$ soit impair pour tout $i$ ;
		\item $\card(E_i\cap E_j)$ soit pair pour tout $i\ne j$.
	\end{itemize}
\end{ex}
\begin{sol}
	On note $M\in\mathcal{M}_{k,n}(\Z)$ la matrice de terme général $[M]_{i,j}=\textbf{1}_{E_i}(j)$. Montrons que $M$ est de rang $k$. On s'intéresse à $MM^T$. En effet :
	$$[MM^T]_{i,j}=\sum_{p=1}^{n}\textbf{1}_{E_i}(p)\textbf{1}_{E_j}(p)=\sum_{p=1}^{n}\textbf{1}_{E_i\cap E_j}(p)=\card(E_i\cap E_j).$$
	En calculant modulo $2$, $\det(MM^T)$ est impair donc non nul, si bien que $MM^T$ est inversible et est par conséquent de rang $k$. \\
	Ensuite, $\rg(MM^T)=\rg(M^T)=\rg(M)$. Passer par les noyaux pour la première égalité. Comme $M$ est de taille $k\times n$, son rang est majoré par ses dimensions, donc $k\leqslant n$. \\
	Une telle famille pour $k=n$ est donnée par $E_i=\{i\}$.
\end{sol}

\begin{ex}
	Soit $n\geqslant1$. Pour tout $\sigma\in\mathfrak{S}_n$, on note $\text{inv}(\sigma)$ le nombre de points de $\llbracket1,n\rrbracket$ invariants par $\sigma$. Calculer : $$\sum_{\sigma\in\mathfrak{S}_n}\frac{\varepsilon(\sigma)}{\text{inv}(\sigma)+1}$$
\end{ex}
\begin{sol}
	Penser à une intégrale ! En effet, si $\sigma\in\mathcal{S}_n$, alors :
	$$\frac{\varepsilon(\sigma)}{\text{inv}(\sigma)+1}=\int_{0}^{1}x^{\text{inv}(\sigma)}\d x.$$ Notons $S_n$ la somme cherchée. Par linéarité de l'intégrale :
	$$S_n=\int_{0}^{1}\sum_{\sigma\in\mathcal{S}_n}\varepsilon(\sigma)x^{\text{inv}(\sigma)}\d x.$$
	La somme fait fortement penser à un déterminant. Mais lequel ? Il s'agit de :
	$$\begin{vmatrix}
		x&&(1)\\
		&\ddots&\\
		(1)&&x
	\end{vmatrix}$$ car le nombre de $x$ choisis correspond, à $\sigma$ fixé, au nombre de points fixes de $\sigma$. Calculons ce déterminant. On somme d'abord toutes les colonnes sur la première puis on utilise la linéarité par rapport à la première colonne.
	\begin{align*}
		\begin{vmatrix}
			x&&(1)\\
			&\ddots&\\
			(1)&&x
		\end{vmatrix}&=\begin{vmatrix}
		x+(n-1)&&&(1)\\
		x+(n-1)&x&&\\
		\vdots&&\ddots&\\
		x+(n-1)&1&&x
		\end{vmatrix}\\
		&=(x+n-1)\begin{vmatrix}
			1&&&(1)\\
			1&x&&\\
			\vdots&&\ddots&\\
			1&(1)&&x
		\end{vmatrix}\\
		&=(x+n-1)\begin{vmatrix}
			1&&&(0)\\
			1&x-1&&\\
			\vdots&&\ddots&\\
			1&(0)&&x-1
		\end{vmatrix}\\
		&=(x+n-1)(x-1)^{n-1}\\
		&=(x-1)^n+n(x-1)^{n-1}.
	\end{align*}
	On obtient finalement :
	$$S_n=\int_{0}^{1}(x-1)^n+n(x-1)^{n-1}\d x=\left[\frac{(x-1)^{n+1}}{n+1}+(x-1)^n\right]_0^1=\frac{(-1)^{n+1}n}{n+1}.$$
\end{sol}

\chapter{Suites et séries numériques}

\section{Problèmes}

\begin{pro}
	Soit $\displaystyle\sum_{n\geqslant 1}a_n$ une série à termes positifs convergente, de somme notée $A$. On pose pour tout $n\in\N\st$ : $u_n=\sqrt[n]{a_1\cdots a_n}$.
	\begin{enumerate}
		\item Démontrer l'inégalité arithmético-géométrique.
		\item Montrer que $\displaystyle\sum_{n\geqslant 1}u_n$ converge et qu'en notant $U$ sa somme, on a $U\leqslant e A$.
		\item Montrer que si une constante $k$ vérifie l'inégalité $U\leqslant k A$ pour toute série à termes positifs convergente $\displaystyle\sum_{n\geqslant 1}a_n$, alors $e\leqslant k$.
	\end{enumerate}
\end{pro}

\begin{pro}
	Convergence et somme de : $$\sum_{n\geqslant 2}\frac{(-1)^n}{n}\left\lfloor\frac{\ln(n)}{\ln(2)}\right\rfloor$$
\end{pro}

\section{Exercices}

\begin{ex}
	Soit $u_0$ un réel strictement positif. On définit par récurrence : $$\forall n\in\N,\ u_{n+1}=u_n^{u_n}$$
	\begin{enumerate}
		\item Etudier la convergence de la suite $(u_n)$.
		\item Désormais et jusqu'à la fin de l'exercice, on suppose $u_0\in\N\setminus\{0,1\}$. Montrer : $$\forall N\in\N,\ \forall n\in\llbracket0,N\rrbracket,\ u_n\mid u_N$$
		\item Montrer que pour tous $N$ et $k$ dans $\N$, on a : $u_{N+k}\geqslant u_N^{k+1}$.
		\item Montrer la convergence de la série de terme général $1/u_n$. On note $S$ sa somme.
		\item Montrer : $$u_N\sum_{n=N+1}^{+\infty}\frac{1}{u_n}\xrightarrow[N\to+\infty]{}0$$
		\item Montrer que $S\notin\Q$.
	\end{enumerate}
\end{ex}

\begin{ex}
	Pour $x\in\R$, étudier la convergence de la suite de terme général $P_n(x)=\displaystyle\prod_{k=1}^{n}\cos\left(\frac{x}{2^k}\right)$.
\end{ex}

\begin{ex}
	On définit la suite $(u_n)_{n\in\N}$ par $u_0=1$ et $$\forall n\in\N,\ u_{n+1}=e^{-u_n}-1$$ Déterminer un équivalent de la suite $(u_n)_{n\in\N}$.
\end{ex}

\begin{ex}
	Déterminer $\displaystyle\underset{n\to+\infty}{\lim}\sum_{k=n}^{+\infty}\frac{e^{n/k}}{k^2}$.
\end{ex}
\begin{sol}
	Fixer $n$ et faire une comparaison série intégrale avec $f:x\mapsto e^{n/x}/x^2$ qui est décroissante. On trouve (après dessin - et une primitive est $F(x)=-e^{n/x}/n$) :
	$$\frac{e}{n}=\int_{n}^{+\infty}f(x)\text{d} x\leqslant\sum_{k=n}^{+\infty}\frac{e^{n/k}}{k^2}\leqslant\int_{n-1}^{+\infty}f(x)\text{d} x=\frac{e^{n/(n-1)}}{n}$$ donc la limite cherchée existe et vaut $e$.
\end{sol}

\chapter{Réduction des endomorphismes}

\section{Problèmes}

\section{Exercices}

\begin{ex}
	Soit $A\in\mathcal{M}_n(\Z/p\Z)$ où $n\in\N\st$ et $p$ est un nombre premier. Montrer que $A$ est diagonalisable si, et seulement si, $A^p=A$.
\end{ex}

\begin{ex}
	Pour $\sigma\in\mathcal{S}_n$, on note $P_\sigma\in\mathcal{M}_n(\C)$ la matrice de permutation associée à $\sigma$, de terme général $\delta_{i,\sigma(j)}$. Soit $(\sigma,\tau)\in\mathcal{S}_n^2$. Montrer que $\sigma$ et $\tau$ sont conjuguées dans le groupe $\mathcal{S}_n$ si, et seulement si, les matrices $P_\sigma$ et $P_\tau$ sont semblables.
\end{ex}

\begin{ex}
	En considérant des sous-groupes dont tous les éléments sont des matrices de symétries, montrer que si $\mathcal{GL}_n(\C)$ et $\mathcal{GL}_m(\C)$ sont isomorphes, alors $n=m$.
\end{ex}
\begin{sol}
	Une matrice de symétrie est diagonalisable sur et ses valeurs propres sont dans ${-1,1}$. Un tel sous-groupe est commutatif et donc toutes les matrices sont simultanément diagonalisables, si bien que le cardinal d'un tel sous-groupe est majoré par $2^n$ (si on est dans $\mathcal{GL}_n(\C)$), et cette borne est atteinte avec les matrices diagonales à coefficients diagonaux dan ${-1,1}$. Si on avait un isomorphisme, un tel sous-groupe s'enverrait sur un sous-groupe similaire mais si $n\ne m$, on a $2^n\ne 2^m$.
\end{sol}

\begin{ex}
	On admet le théorème de réduction de Jordan. Montrer que toute matrice inversible à coefficients complexes admet une racine carrée. Démontrer le théorème de réduction de Jordan.
\end{ex}
\begin{sol}
	Commencer par le cas d'une racine carrée de $I+N$ avec $N$ nilpotente. Utiliser le DL de $\sqrt{1+x}$ pour obtenir le bon candidat. Le cas général s'en déduit facilement par homothétie. Pour Jordan, comme d'habitude...
\end{sol}

\chapter{Espaces euclidiens}

\section{Problèmes}

\begin{pro}[Autour du théorème de Schur-Horn]
	
\end{pro}

\begin{pro}[Autour des égalités de Courant-Fischer]
	On considère $u\in S(E)$ avec $E$ euclidien de dimension $n$. On note $\lambda_1\leqslant\ldots\leqslant\lambda_n$ ses valeurs propres. On note $\mathcal{F}_k$ l'ensemble des sous-espaces vectoriels de $E$ de dimension $k$ et $S$ la sphère unité de $E$.
	\begin{enumerate}
		\item Montrer que : $$\lambda_k=\inf_{F\in\mathcal{F}_k}\sup_{x\in S\cap F}(x|u(x))$$ puis montrer que $$\lambda_{n+1-k}=\sup_{F\in\mathcal{F}_k}\inf_{x\in S\cap F}(x|u(x))$$
		\item Soit $A\in\mathcal{S}_n(\R)$ avec $n\geqslant 2$. On note $B$ la matrice extraite de $A$ en lui supprimant sa dernière ligne et sa dernière colonne. On note $\lambda_1\leqslant\ldots\leqslant\lambda_n$ les valeurs propres de $A$ et $\mu_1\leqslant\ldots\leqslant\mu_{n-1}$ celles de $B$. Montrer : $$\lambda_1\leqslant\mu_1\leqslant\lambda_2\leqslant\ldots\leqslant\mu_{n-1}\leqslant\lambda_n$$
		\item Pour $A\in\mathcal{S}_n(\R)$, on note $\lambda_1(A)\geqslant\ldots\geqslant\lambda_n(A)$ les valeurs propres de $A$, comptées avec multiplicité. Montrer : $$\forall (A,B)\in\mathcal{S}_n(\R)^2,\ \lambda_{i+j-1}(A+B)\leqslant\lambda_i(A)+\lambda_j(B)$$
	\end{enumerate}
\end{pro}

\begin{pro}[Autour des racines carrées]
	On présente ici la racine carré d'un endomorphisme symétrique positif et quelques applications.
	\begin{enumerate}
		\item Soit $E$ un espace euclidien ainsi que $u\in\mathcal{S}^+(E)$. Montrer qu'il existe un unique endomorphisme $v\in\mathcal{S}^+(E)$ tel que $u=v^2$. Montrer que $v$ est un polynôme en $u$.
		\item Soit $n\in\N\st$ et $M\in\mathcal{GL}_n(\R)$. Montrer qu'il existe un unique couple $(O,S)\in\mathcal{O}_n(\R)\times\mathcal{S}_n^{++}(\R)$ tel que $M=OS$. Montrer que l'on conserve l'existence si l'on suppose seulement $M\in\mathcal{M}_n(\R)$. Montrer que l'application qui à $(O,S)\in\mathcal{O}_n(\R)\times\mathcal{S}_n^{++}(\R)$ associe $OS$ est un homéomorphisme de $\mathcal{O}_n(\R)\times\mathcal{S}_n^{++}(\R)$ sur $\mathcal{GL}_n(\R)$.
		\item Montrer que pour toute matrice diagonalisable à coefficients réels $D$, il existe $S\in\mathcal{S}_n^{++}(\R)$ et $T\in\mathcal{S}_n(\R)$ telles que $D=ST$.
		\item Soit $A\in\mathcal{S}_n^{++}(\R)$ et $B\in\mathcal{S}_n(\R)$. Montrer qu'il existe $P\in\mathcal{GL}_n(\R)$ telle que $P^TAP=I_n$ et $P^TBP$ soit diagonale.
	\end{enumerate}
\end{pro}

\section{Exercices}

\begin{ex}
	On projette orthogonalement un cube de l'espace de côté $a$ sur un plan. On note $x$, $y$ et $z$ les longueurs des arêtes du projeté. Montrer que $x^2+y^2+z^2=2a^2$.
\end{ex}
\begin{sol}
	Faire un dessin. Remarquer que dès que l'on dessin un cube en perspective cavalière, il est déjà projeté dans un plan. Considérer ensuite deux bases orthonormées de l'espace : une base liée au cube, et une base du plan de projection complétée en une base de l'espace. Écrire la matrice de passage entre ces deux bases (qui est orthogonale). Conclure en utilisant une relation sur les lignes. 
\end{sol}

\begin{ex}
	Soit $E$ un espace euclidien de dimension $n$. On suppose qu'il existe $(x_1,\ldots,x_p)\in E^{p}$ telle que : $$\forall (i,j)\in\llbracket1,p\rrbracket^2,\ i\ne j\implies (x_i\mid x_j)<0$$ Montrer que $p\leqslant n+1$.
\end{ex}
\begin{sol}
	Par récurrence comme des sagouins mais je n'ai jamais essayé. Sinon, on peut montrer que $(x_1,\ldots,x_{p-1})$ est libre : séparer une relation de liaison entre les scalaires positifs et les scalaires strictement négatifs, faire le produit scalaire avec le vecteur correspondant aux scalaires positifs. On trouve une norme négative, donc les vecteurs avec les scalaires positifs et strictement négatifs sont tous les deux nuls. En effectuant le produit scalaire de ces deux vecteurs avec $x_p$ et par un argument de signe, on conclut.
\end{sol}

\begin{ex}
	Soit $E$ un espace euclidien. Montrer que le groupe $\mathcal{O}(E)$ des isométries vectorielles est engendré par l'ensemble des réflexions.
\end{ex}
\begin{sol}
	On peut le faire en utilisant le théorème de réduction des isométries vectorielles en montrant que les blocs de rotation s'obtiennent comme le produit de deux réflexions. \\
	Sinon, on peut raisonner par récurrence descendante sur la dimension de l'espace propre associé à la valeur propre $1$.\\
	On peut aussi le faire par récurrence sur la dimension de $E$ en projetant sur : un vecteur invariant s'il existe, ou bien sur un hyperplan médiateur s'il n'y a pas de vecteur invariant.
\end{sol}

\begin{ex}
	Soit $n\in\N\st$ et $M\in\mathcal{M}_n(\R)$. On suppose qu'il existe un entier $k\geqslant 2$ tel que $M^k=M^T$.
	\begin{enumerate}
		\item Montrer que $M^TM$ est la matrice d'une projection orthogonale.
		\item Montrer que $M$ est $\C$-diagonalisable.
	\end{enumerate}
\end{ex}

\begin{ex}
	Soit $(A,B)\in\mathcal{S}_n(\R)^2$ vérifiant : $$\forall X\in\R^n,\ 0\leqslant X^TAX\leqslant X^TBX$$ Montrer : $\det(A)\leqslant\det(B)$.
\end{ex}

\begin{ex}
	Soit $S\in\mathcal{S}_n^{++}(\R)$ et $A\in\mathcal{A}_n(\R)$. Montrer que $SA$ est diagonalisable sur $\C$ et à spectre inclus dans $i\R$.
\end{ex}
\begin{sol}
	Notons $R$ la racine carrée de $S$ (définie positive). Alors $SA=R(RAR)R\inv$ donc $SA$ est semblable à $RAR$. Or, cette dernière est antisymétrique, puis on démontre les résultats classiques sur les matrices antisymétriques.
\end{sol}

\chapter{Topologie}

\section{Problèmes}

\begin{pro}[Autour de la propriété de Borel-Lebesgue]
	On se donne $K$ un compact.
	\begin{enumerate}
		\item  Précompacité : montrer que pour tout $\varepsilon>0$, il existe une partie finie $C$ de $K$ telle que $$K\subset\bigcup_{x\in C}B(x,\varepsilon)$$
		\item Propriété de Borel-Lebesgue :
		\begin{enumerate}[label=\alph*.]
			\item Soit $(O_i)_{i\in I}$ une famille d'ouverts de $K$ recouvrant $K$. Montrer qu'il existe $\varepsilon>0$ tel que $$\forall x\in K,\ \exists i\in I\ : \ B_o(x,\varepsilon)\subset O_i$$
			\item Montrer que pour toute famille $(O_i)_{i\in I}$ d'ouverts de $A$ recouvrant $K$, il existe une partie finie $J$ de $I$ telle que $(O_j)_{j\in J}$ recouvre $K$. On pourra utiliser la précompacité de $K$.
			\item Montrer que la propriété de la deuxième question est suffisante pour que $K$ soit compacte. On pourra utiliser le fait que l'ensemble des valeurs d'adhérence d'une suite $(u_n)$ est $\displaystyle\bigcap_{n\in\N}\overline{\{u_k\mid k\geqslant n\}}$
		\end{enumerate}
		\item Applications :
		\begin{enumerate}[label=\alph*.]
			\item Montrer que les idéaux maximaux de $\mathcal{C}^0([0,1],\R)$ sont de la forme : $$I_x=\left\{f\in\mathcal{C}^0([0,1],\R)\ \mid\ f(x)=0\right\}$$
			\item Soit $E$ un espace métrique compact et $f:E\to\R$ une fonction localement bornée, c'est-à-dire que pour tout $x\in E$, il existe un voisinage $V_x$ de $x$ tel que pour tout $y\in V_y$, on ait $\lvert f(y)\rvert\leqslant M_x$. Montrer que $f$ est bornée.
			\item On admet le théorème de Kakutani : pour tout convexe compact non vide $K$ d'un espace vectoriel normé $E$ et toute application affine continue $u$ qui stabilise $K$, $u$ possède un point fixe. Montrer qu'une famille $(u_i)_{i\in I}$ d'applications affines qui commutent deux à deux et stabilisent toutes un compact convexe non vide $K$ possèdent un point fixe commun.
		\end{enumerate}
	\end{enumerate}
\end{pro}

\begin{pro}[Autour du théorème de Baire]
	Soit $A$ une partie localement compacte d'un espace vectoriel normé (\textit{ie} dans laquelle tout point possède un voisinage compact).
	\begin{enumerate}
		\item Soit $(O_n)$ une suite d'ouverts de $A$, tous denses dans $A$. Montrer que $$\Omega=\bigcap_{n\in\N} O_n$$ est dense dans $A$. Est-il nécessairement ouvert ?
		\item En déduire que si une suite $(F_n)$ de fermés de $A$ a une réunion d'intérieur non vide, alors l'un des $F_n$ est d'intérieur non vide.
	\end{enumerate}
\end{pro}

\begin{pro}[Connexités par arcs]
	\begin{enumerate}
		\item Montrer que l'ensemble $\mathcal{GL}_n(\C)$ est connexe par arcs.
		\item Montrer que $\mathcal{GL}_n(\R)$ n'est pas connexe par arcs et déterminer ses composantes connexes par arcs.
		\item Déterminer les composantes connexes par arcs de l'ensemble des matrices de projecteurs pour $\K=\C$ puis pour $\K=\R$.
	\end{enumerate}
\end{pro}

\section{Exercices}

\begin{ex}[Point fixe de Kakutani]
	Soit $K$ un compact convexe non vide et $u$ une application affine continue qui stabilise $K$. Montrer que $u$ admet un point fixe.
\end{ex}

\begin{ex}
	Déterminer les sous-groupes fermés de $(\C\st,\times)$.
\end{ex}

\begin{ex}
	Montrer que pour tout polynôme $P\in\R[X]$, il existe une norme sur $\R[X]$ qui fait converger la suite $(X^k)_{k\in\N}$ vers $P$.
\end{ex}

\begin{ex}
	Montrer qu'une matrice $A\in\mathcal{M}_n(\K)$ est nilpotente si, et seulement si, la matrice nulle est adhérente à la classe de similitude de $A$.
\end{ex}

\begin{ex}
	Montrer qu'il n'existe pas de norme sur $\mathcal{C}^{\infty}([0,1],\R)$ pour laquelle la dérivation soit continue.
\end{ex}

\begin{ex}
	Soit $C$ une partie convexe compacte non vide et $f:C\to C$ 1-lipschitzienne. Montrer que $f$ admet au moins un point fixe. On pourra commencer par le cas où $f$ est $k$-lipschitzienne avec $k<1$.
\end{ex}

\begin{ex}
	Soit $A$ une partie non triviale d'un espace vectoriel normé $E$. Montrer que la frontière de $A$ est non vide.
\end{ex}

\begin{ex}
	Soit $G$ un sous-groupe additif de $\R^n$. On suppose que $G$ est fermé et que pour tout $y\in G$ et pour tout $\varepsilon>0$, il existe $x\in G$ tel que $0<\lVert x-y\rVert<\varepsilon$. Montrer que $G$ contient une demi-droite.
\end{ex}

\begin{ex}
	Soit $(E,d)$ un espace métrique et $K$ une partie compacte de $E$. On considère $f$ une fonction de $K$ dans $K$ telle que : $$\forall (x,y)\in K^2,\ d(f(x),f(y))\geqslant d(x,y)$$
	Montrer que $f(K)$ est dense dans $K$ et en déduire que $f$ est un homéomorphisme.
\end{ex}

\begin{ex}[Lemme d'Auerbach]
	Soit $(E,\lVert\cdot\rVert)$ un espace vectoriel normé de dimension finie. Montrer qu'il existe une base de $E$ unitaire dont la base duale est unitaire (pour la norme subordonnée).
\end{ex}

\chapter{Fonctions vectorielles d'une variable réelle}

\section{Problèmes}

\section{Exercices}

\chapter{Suites et séries de fonctions}

\section{Problèmes}

\begin{pro}[Théorème de Weierstrass trigonométrique]
	Soit $f$ une fonction continue de $\R$ dans $\C$ $2\pi$-périodique.
	\begin{enumerate}
		\item En utilisant le théorème de Weierstrass, montrer que pour toute fonction $f\in\mathcal{C}^0(\R,\C)$ paire et $2\pi$-périodique, il existe une suite de polynômes $(P_n)_{n\in\N}$ telle que la suite de fonctions $(t\mapsto P_n(\cos(t)))_{n\in\N}$ converge uniformément vers $f$.
		\item En déduire qu'il existe une suite de polynômes trigonométriques qui converge uniformément sur $\R$ vers $f\times\sin^2$.
		\item En déduire qu'il existe une suite de polynômes trigonométriques qui converge uniformément vers $f$ sur $\R$.
	\end{enumerate}
\end{pro}
\begin{sol}
	\begin{enumerate}
		\item On considère $g=f\circ\arccos$ définie sur $[-1,1]$ et on lui applique le théorème de Weierstrass pour obtenir $(P_n)$. Pour montrer que la convergence est alors uniforme sur $\R$, utiliser la parité et la $2\pi$-périodicité de $f$ et de $\cos$.
		\item Utiliser le fait que $f$ s'écrit sous la forme $p+i$ avec $p$ paire et $i$ impaire (où $p$ et $i$ sont alors nécessairement $2\pi$-périodiques car elle s'expriment naturellement à partir de $f$). On remarque alors que : $$f\times\sin^2=p\times\sin^2+\sin\times(i\times\sin)$$ où $p\times\sin^2$ et $i\times\sin$ sont paires et $2\pi$ périodiques. On les approche alors par des polynômes trigonométriques $(P_n\circ\cos)$ et $(Q_n\circ\cos)$ grâce à l'exercice précédent, puis on montre ensuite que $(P_n\circ\cos+\sin\times Q_n\circ\cos)$ est une suite de polynômes trigonométriques qui approche uniformément $f\times\sin^2$.
		\item En considérant les transformations $x\mapsto x+\displaystyle\frac{\pi}{2}$ et $x\mapsto x-\displaystyle\frac{\pi}{2}$ (en appliquer une à $f$), on se ramène à la question précédente pour montrer que $f\times\cos^2$ peut être approchée uniformément par une suite de polynômes trigonométriques. On conclut alors avec $\cos^2+\sin^2=1$.
	\end{enumerate}
\end{sol}

\begin{pro}
	\begin{enumerate}
		\item On considère une suite décroissante de fonctions $(f_n)$ sur un compact $K$ à valeurs dans $\R$ qui converge simplement vers $0$. Montrer que la convergence est uniforme.
		\item On considère une suite de fonctions $(f_n)$ définies sur $[a,b]$ à valeurs réelles qui sont toutes croissantes et qui converge simplement vers $f$ continue. Montrer que la convergence est uniforme.
		\item Montrer que si une suite de fonctions convexes sur un intervalle ouvert $I$ converge simplement sur $I$, alors la convergence est uniforme sur tout segment de $I$.
	\end{enumerate}
\end{pro}
\begin{sol}
	\begin{enumerate}
		\item Déjà, remarquons que toutes les $f_n$ sont par conséquent positives. Il est possible de démontrer le résultat à la main en prenant les \textit{maxima} sur le segment, mais il y a plus élégant. Soit $\varepsilon>0$. On considère $$K_n:=\left\{x\in K \ : \ f_n(x)\geqslant\varepsilon\right\}$$ puis on raisonne par l'absurde. On suppose que tous ces compacts (fermés dans un compact) sont non vides. Alors, puisque $(K_n)$ forme une suite décroissante de compacts, d'après le théorème des compacts emboîtés, on peut prendre $x$ dans l'intersection de tous ces compacts. Mais alors il ne peut y avoir convergence simple en $x$. Cela est absurde. Ainsi, il existe un $K_{n_0}$ vide, et par décroissance, tous les suivants le sont. On a donc bien convergence uniforme.
		\item Un passage à la limite montre que $f$ est croissante. On fixe $\varepsilon>0$ Et par croissance et uniforme continuité de $f$, on peut fixer $c_0<\ldots<c_n$ une subdivision de $[0,1]$ telle que $$\forall i,\ f(c_{i+1})-f(c_i)\leqslant\frac{\varepsilon}{2}$$ Ensuite, si $x\in[0,1]$, on prend $i$ tel que $c_i\leqslant x\leqslant c_{i+1}$. Par croissance de $f$ et de $f_n$ on a alors : $$\begin{cases}
			f_n(x)-f(x)\leqslant f_n(c_{i+1})-f(c_i)\leqslant\frac{\varepsilon}{2}+f_n(c_{i+1})-f(c_{i+1})\\
			f_n(x)-f(x)\geqslant f_n(c_i)-f(c_{i+1})\geqslant f_n(c_i)-f(c_i)-\frac{\varepsilon}{2}
		\end{cases}$$ Or, il y a un nombre fini de $c_i$ donc on peut prend $n_0$ tel que $$\forall i,\ \forall n\geqslant n_0,\ \lvert f_n(c_i)-f(c_i)\rvert\leqslant\frac{\varepsilon}{2}$$ On a donc : $$\forall x,\forall n\geqslant n_0,\ \lvert f_n(x)-f(x)\rvert\leqslant\varepsilon$$ ce qui est bien la convergence uniforme souhaitée.
		\item On note $(f_n)$ la suite de fonctions et $f$ la limite simple. On prend $[b,c]\in I$ puis $(a,d)\in I^2$ tel que $a<b<c<d$. Soit $x\ne y\in[b,c]$. La convexité fournit : $$\frac{f_n(b)-f_n(a)}{b-a}\leqslant\frac{f_n(y)-f_n(x)}{y-x}\leqslant\frac{f_n(d)-f_n(c)}{d-c}$$ Or, les deux suite extrémales admettent des limites indépendantes de $x$ et $y$. Donc on en déduit que $$\forall n,\ \forall x\ne y,\ \left\lvert \frac{f_n(y)-f_n(x)}{y-x}\right\rvert\leqslant M$$ où $M$ ne dépend pas ni de $x$, ni de $y$, ni de $n$. Ensuite, il s'agit d'un exercice classique où toutes les fonctions sont lipschitziennes de même rapport.
	\end{enumerate}
\end{sol}

\section{Exercices}

\begin{ex}
	Soit $a\in\displaystyle\left]0,\frac{1}{2}\right[$.
	\begin{enumerate}
		\item Montrer que la suite de fonctions $\displaystyle\left(x\mapsto(1+(2x-1)^n)/2\right)_{n\in\N}$ converge uniformément vers $x\mapsto\displaystyle\frac12$ sur $[a,1-a]$.
		\item En déduire que toute fonction continue de $[a,1-a]$ dans $\R$ est limite uniforme d'une suite de fonction polynomiales à coefficients entiers.
		\item Toute fonction continue de $[0,1]$ dans $\R$ est-elle limite uniforme d'une suite de fonctions polynomiales à coefficients entiers ?
	\end{enumerate}
\end{ex}
\begin{sol}
	\begin{enumerate}
		\item La limite simple est rapide avec une suite géométrique. Pour la convergence uniforme, majorer la raison en module par $0<1-2a<1$ pour l'uniformité.
		\item Vérifier que la suite de fonctions précédente est polynomiale à coefficients entiers (avec un binôme de Newton cela se fait bien). \\ Ainsi, $1/2$ est limite uniforme d'une suite de fonctions polynomiales à coefficients entiers. \\ Par produit $p$ fois, on en déduit que tous les $1/2^p$ sont limite uniforme d'une suite de fonctions polynomiales à coefficients entiers. \\  Avec les dyadiques et par somme, on en déduit que toute fonction constante est limite uniforme de fonctions polynomiales à coefficients entiers. \\ On en déduit, par produit, avec des puissances de $x$, que toute fonction polynomiale sur $[-a,1-a]$ est limite uniforme de fonctions à coefficients entiers. \\ Puis on conclut avec Weierstrass et le fait que la limite uniforme possède une sorte d'associativité.
		\item Au vu de la question, je dirais non. Il doit y avoir un problème lorsque l'on s'approche de $0$ et $1$ au vu de la construction qui précède. \\  Après réflexion, prendre une fonction qui vaut $1/2$ en $0$. Comme $\Z$ est fermé dans $\R$, il va y avoir un problème avec la convergence uniforme qui entrainerait le fait qu'une suite d'entiers pourrait converger vers $1/2$.
	\end{enumerate}
\end{sol}

\chapter{Intégrales généralisées}

\begin{ex}
	Pour $n\in\N$ et $x\in\R$, on pose : $H_n(x)=\displaystyle\frac{1}{n!}\int_{-x}^{x}(x^2-t^2)^n\cos(t)\ dt$
	\begin{enumerate}
		\item Montrer que pour tout $n\in\N$, il existe des polynômes $C_n$ et $S_n$ à coefficients entiers de degrés inférieurs ou égaux à $n$ tels que pour tout $x\in\R$, on ait : $H_n(x)=C_n(x)\cos(x)+S_n(x)\sin(x)$.
		\item Soit $\theta\in\Q\st$. Montrer que $\tan(\theta)$ est irrationnel.
	\end{enumerate}
\end{ex}

\begin{ex}
	Calculer $\displaystyle\int_{0}^{\pi/2}\ln(\sin(x))dx$.
\end{ex}

\begin{ex}
	Calculer $\displaystyle\int_{0}^{+\infty}\frac{\arctan(x)^2}{x^2}\text{d}x$.
\end{ex}
\begin{sol}
	L'existence est simple. Pour le calcul, on peut faire une IPP en primitivant $x^{-2}$, puisque on effectue le changement de variable $x=\tan(u)$. Cela permet de se ramener à l'intégrale entre $0$ et $\pi/2$ de $\ln(\sin(u))$, que l'on calcule en faisant un changement de variable puis en sommant avec $\ln(\cos)$ puis en refaisant un changement de variable. On trouve pour celle-ci $-\pi\ln(2)/2$, et finalement $\pi\ln(2)$ pour l'intégrale de départ.
\end{sol}

\chapter{Intégrales à paramètres}

\section{Problèmes}

\begin{pro}
	Pour $(x,y)\in\R_+\st\times\R$, on pose $$G(x,y)=\int_{0}^{+\infty}\frac{e^{-xt}e^{iyt}}{\sqrt{t}}\ dt$$
	\begin{enumerate}
		\item Montrer que $G$ est continue sur $\R_+\st\times\R$.
		\item Calculer $G(1,0)$.
		\item Exprimer $G(x,y)$ à l'aide de $G(1,z)$ pour un $z$ bien choisi.
		\item Montrer que $H:z\mapsto G(1,z)$ est dérivable sur $\R$ et que $$\forall z\in\R,\ H'(z)=\frac{i}{2(1-iz)}H(z)$$
		\item En déduire la valeur de $H(z)$ pour tout $z\in\R$ puis celle de $G(x,y)$ pour tout $(x,y)\in\R_+\st\times\R$.
		\item Montrer que $G(0,y)$ a un sens pour tout $y>0$.
		\item Montrer que $\displaystyle\lim_{x\to 0^+}G(x,1)=G(0,1)$. En déduire la convergence et la valeur de $$I=\int_{0}^{+\infty}\cos(t^2)\ dt \ \ \text{et} \ \ J=\int_{0}^{+\infty}\sin(t^2)\ dt$$
	\end{enumerate}
\end{pro}

\begin{pro}
	Soit $f$ une fonction continue et à valeurs strictement positives sur $[0,a]$ avec $a>0$. on suppose qu'il existe $(\alpha,p)\in(\R_+\st)^2$ tel que $f(t)=1-\alpha t^p+o(t^p)$ au voisinage de $0$ et $\forall t\in]0,a],\ f(t)<1$. On se propose de trouver un équivalent de $$I(x)=\int_{0}^{a}f(t)^x\ dt$$ lorsque $x$ tend vers $+\infty$.
	\begin{enumerate}
		\item Montrer qu'il existe $\beta>0$ tel que $$\forall t\in[0,a],\ f(t)\leqslant\exp(-\beta t^p)$$
		\item Déterminer la limite, quand $x$ tend vers $+\infty$, de $$g(x)=\int_{0}^{ax^{1/p}}f\left(ux^{-1/p}\right)^x\ dt$$
		\item Conclure.
		\item Application : convergence, limite et équivalent lorsque $n$ tend vers l'infini de la suite $(I_n)$ définie par $$I_n=\int_{0}^{+\infty}\frac{dt}{(1+t^3)^n}$$
	\end{enumerate}
\end{pro}

\begin{pro}
	Si $f$ et $g$ sont des fonctions continues par morceaux de $\R$ dans $\R$, on appelle produit de convolution de $f$ et $g$, $f\star g$, la fonction, si elle existe, définie par $$f\star g:x\mapsto\displaystyle\int_\R f(x-t)g(t)\ dt$$
	\begin{enumerate}
		\item Montrer que si $f$ et $g$ sont à support compact, alors $f\star g$ est aussi à support compact.
		\item Démontrer que le produit de convolution est commutatif. On admet qu'il est associatif.
		\item Montrer que si $f$ est bornée et $g$ est intégrable, alors $f\star g$ existe et est bornée par $\lVert f\rVert_\infty\lVert g\rVert_1$.
		\item Pour $a>0$, on note $\varphi_a$ la fonction $t\mapsto\exp(-a^2t^2)$. Montrer que $$\varphi_a\star\varphi_b=\sqrt{\frac{\pi}{a^2+b^2}}\varphi_{\frac{ab}{\sqrt{a^2+b^2}}}$$
		\item Soit $f$ continue bornée sur $\R$. Soit $(\varphi_n)$ une suite de fonctions continues positives intégrables telle que $\displaystyle\int_\R\varphi_n(t)\ dt\xrightarrow[n\to+\infty]{}1$ et, pour tout $\delta>0$, $\displaystyle\int_{\lvert t\rvert\geqslant \delta}\varphi_n(t)\ dt\xrightarrow[n\to+\infty]{}0$. Montrer que la suite de fonctions $(f\star\varphi_n)$ converge simplement vers $f$, et uniformément sur tout segment.
		\item Pour $n\in\N$, on pose $\displaystyle g_n(x)=\frac{1}{c_n}(1-x^2)^n$ si $x\in[-1,1]$ avec $c_n=\displaystyle\int_{-1}^{1}(1-t^2)^n\ dt$. Soit $f:[-1,1]\to\R$ nulle en dehors de $\displaystyle\left[-\frac{1}{2},\frac{1}{2}\right]$. Montrer que la restriction de $f\star g_n$ à $\displaystyle\left[-\frac{1}{4},\frac{1}{4}\right]$ est polynomiale.Montrer que la suite de fonctions $(f\star g_n)$ converge uniformément vers $f$ sur $\displaystyle\left[-\frac{1}{4},\frac{1}{4}\right]$ et en déduire une preuve du théorème de Weierstrass.
	\end{enumerate}
\end{pro}

\section{Exercices}

\begin{ex}[Démonstration du théorème de D'Alembert-Gauss]
	Soit $f$ une fonction de $\R$ dans $\C\st$, $2\pi$-périodique et de classe $\mathcal{C}^1$.
	\begin{enumerate}
		\item Montrer : $$\frac{1}{2i\pi}\int_{0}^{2\pi}\frac{f'(t)}{f(t)}dt\in\Z$$
		\item En déduire le théorème de D'Alembert-Gauss.
		\item Calculer la valeur de cette intégrale pour $P\in\C[X]$.
	\end{enumerate}
\end{ex}

\begin{ex}
	Soit $a$ et $b$ deux réels strictement supérieurs à $1$. Calculer : $$\int_{0}^{\pi}\ln\left(\frac{b-\cos(x)}{a-\cos(x)}\right)dx$$
\end{ex}
\begin{sol}
	Interpréter ceci comme $I(b)-I(a)$. On dérive $x\mapsto I(x)$, on calcule sa dérivée avec les règles de Bioche, puis on intègre. Inutile de déterminer la constante puisque ce que l'on cherche est la différence entre deux valeurs de $I$.
\end{sol}

\begin{ex}[Formule sommatoire de Poisson]
	On note $\mathcal{S}$ l'ensemble des fonctions de $\R$ dans $\C$ de classe $\mathcal{C}^\infty$ telles que pour tous entiers naturels $n$ et $m$, on ait $x^mf^{(n)}(x)\xrightarrow[x\to\pm\infty]{}0$. Soit $f\in\mathcal{S}$. On pose : $$\forall x\in \R,\ \varphi(x)=\sum_{k\in\Z}f(x+k)$$
	\begin{enumerate}
		\item Bonne définition de $\varphi$ ? Régularité de $\varphi$ ? Une autre remarque ?
		\item Pour $n\in\Z$, calculer : $$\int_{0}^{1}\varphi(t)e^{-2i\pi nt}dt$$
		\item On admet qu'une fonction $\psi$ de classe $\mathcal{C}^\infty$ et $1$-périodique s'écrit sous la forme : $$\psi(x)=\sum_{n\in\Z}c_ne^{2i\pi nx}$$ On pose : $$\forall n\in\Z,\ \hat{f}(n)=\int_{\R}f(t)e^{-2i\pi nt}dt$$ Montrer la formule sommatoire de Poisson : $$\sum_{n\in\Z}f(n)=\sum_{n\in\Z}\hat{f}(n)$$
	\end{enumerate}
\end{ex}

\begin{ex}
	Soit $f:\R\to\C$ continue et bornée. Démontrer l'existence de $$I(\varepsilon)=\int_{-\infty}^{+\infty}\frac{\varepsilon}{x^2+\varepsilon^2}f(x)\d x$$ et étudier son comportement quand $\varepsilon\xrightarrow{}0$.
\end{ex}

\chapter{Probabilités}

\section{Problèmes}

\section{Exercices}

\begin{ex}
	On considère le déplacement d'une fourmi sur les sommets d'un carré $ABCD$, avec le segment $[AB]$ horizontal. A l'étape $0$, la fourmi est en $A$. A chaque étape, la fourmi change nécessairement de sommet. Un déplacement horizontal à une probabilité $p$, un déplacement vertical une probabilité $q$ et un déplacement diagonal une probabilité $r$ (avec $p+q+r=1$ donc). Donner, pour chaque sommet, la probabilité que la fourmi s'y trouve à l'étape $n$. Limite quand $n$ tend vers $+\infty$ ?
\end{ex}
\begin{ex}
	On note $a_n$, $b_n$, $c_n$ et $d_n$ ces probabilités. On a donc $(a_0,b_0,c_0,d_0)=(1,0,0,0)$ et :
	$$\forall n\in\N, \begin{pmatrix}
		a_{n+1}\\
		b_{n+1}\\
		c_{n+1}\\
		d_{n+1}
	\end{pmatrix}=\underbrace{\begin{pmatrix}
		0&p&r&q\\
		p&0&q&r\\
		r&q&0&p\\
		q&r&p&0
	\end{pmatrix}}_{:=M}\begin{pmatrix}
		a_n\\
		b_n\\
		c_n\\
		d_n
	\end{pmatrix}.$$
	Remarquer que $M$ est symétrique réelle !!! On peut aussi voir que $M$ est la combinaison linéaire de trois matrices de symétrie qui commutent deux à deux... (Et $M$ est bistochastique). On cherche donc à calculer les puissances de $M$. Plusieurs pistes :
	\begin{itemize}
		\item multinôme de Newton comme les matrices de symétrie commutent deux à deux (pas beau du tout) ;
		\item diagonaliser (pas beau non plus, mais plus faisable).
	\end{itemize}
	Je pars sur la deuxième piste. On voit assez aisément que $1$ est valeur propre avec, par exemple, le vecteur propre $(1\ 1\ 1\ 1)^T$.
\end{ex}

\chapter{Équations différentielles}

\section{Problèmes}

\section{Exercices}

\begin{ex}
	Soit $f\in\mathcal{C}^\infty(\R,\R)$ telle que l'espace vectoriel $\text{Vect}(x\mapsto f(x+a))_{a\in\R}$ soit de dimension finie. Montrer que $f$ est solution d'une équation différentielle linéaire à coefficients constants.
\end{ex}

\chapter{Calcul différentiel}

\section{Problèmes}

\section{Exercices}

\begin{ex}
	Soit $A\in\mathcal{M}_n(\R)$. On considère l'application :
	$$\begin{array}{lrcl}
		\varphi_A:&\mathcal{GL}_n(\R)&\longrightarrow&\mathcal{M}_n(\R)\\
		&M&\longmapsto&MAM\inv.
	\end{array}$$
	\begin{enumerate}
		\item Montrer que $\varphi_A$ est différentiable et que le rang de $\d \varphi_A(M)$ ne dépend pas de $M\in\mathcal{GL}_n(\R)$.
		\item Calculer ce rang en fonction des éléments propres de $A$ lorsque $A$ est diagonalisable.	
	\end{enumerate}
\end{ex}
\begin{sol}
	\begin{enumerate}
		\item Se "souvenir" que la fonction inverse est différentiable et que sa différentielle au point $M$ est $H\mapsto -M\inv HM\inv$. Cela donne par opérations la différentiabilité de $\varphi_A$ et, après calcul, on trouve :
		$$\d\varphi_A(M)\cdot H=M(M\inv HA-AM\inv H)M\inv=M([M\inv H, A])M\inv$$ où $[X,Y]=XY-YX$. On peut donc voir $\d\varphi_A(M)$ comme la composée de trois applications linéaires : $X\mapsto MXM\inv$, $X\mapsto [X,A]$ et $X\mapsto M\inv X$. La première et la dernière sont des isomorphismes, donc n'affectent pas le rang. Ainsi, le rang de la différentielle $\d\varphi_A(M)$ est en fait le rang de l'application $X\mapsto [X,A]$, qui ne dépend pas de $M$.
		\item On cherche donc le rang de $X\mapsto [X,A]$, ce qui revient, par le théorème du rang, à chercher la dimension du noyau de $X\mapsto[X,A]$, soit la dimension du commutant de $A$. Or, c'est un exercice classique de réduction que de montrer que pour une matrice diagonalisable, on a :
		$$\dim(\mathcal{C}(A))=\sum_{\lambda\in\text{Sp}(A)}m(\lambda)^2.$$En effet, on montrer qu'un endomorphisme commute avec l'endomorphisme canoniquement associé à $A$ si, et seulement si, ses induits sur les sous-espaces propres commutent avec les induits de $A$.\\
		Finalement :
		$$\rg(\d\varphi_A(M))=n^2-\sum_{\lambda\in\text{Sp}(A)}m(\lambda)^2.$$ 
	\end{enumerate}
\end{sol}

\chapter{Algèbre générale}

\section{Problèmes}

\begin{pro}[Théorème de Wedderburn]
	Pour $n\in\N\st$, on note $\Phi_n$ le $n$-ième polynôme cyclotomique usuel. On se donne $\K$ un corps fini \textit{a priori} non commutatif (on parle de corps gauche). On note $Z$ son centre. On admet que la théorie des espaces vectoriels est inchangée si le corps de base est un corps gauche. Le but est de montrer que $\K$ est en fait commutatif.
	\begin{enumerate}
		\item Soit $n\in\N\st$. Montrer : $X^n-1=\displaystyle\prod_{d\mid n}\Phi_d$. En déduire que les $\Phi_n$ sont à coefficients entiers.
		%\item Soit $p$ un nombre premier et $n\in\N\st$. Montrer : $\Phi_{p^n}=\displaystyle\sum_{k=0}^{p-1}X^{kp^{n-1}}$.
		%\item Expliciter le terme constant de $\Phi_n$ en fonction de $n\in\N\st$.
		\item Soit $d$ un diviseur strict de $n\in\N\st$. Montrer que $X^d-1$ divise $X^n-1$ dans $\Z[X]$, puis que $\Phi_n$ divise $(X^n-1)/(X^d-1)$ dans $\Z[X]$.
		\item Soit $\mathbb{M}$ un corps (gauche) et $\L$ un sous-corps (gauche) de $\mathbb{M}$. Expliquer brièvement pourquoi $\mathbb{M}$ peut être muni d'une structure de $\L$-espace vectoriel et, dans le cas où $\mathbb{M}$ est fini, donner une relation entre les cardinaux de $\mathbb{M}$ et $\L$, et la dimension de $\mathbb{M}$ en tant que $\L$-espace vectoriel.
		\item Montrer que $Z$ est un corps commutatif. On note $q$ le cardinal de $Z$ et $n$ la dimension de $\K$ en tant que $Z$-espace vectoriel. Exprimer alors le cardinal de $\K$ en fonction de $q$ et $n$.
		\item Soit $a\in\K\st$. On note $\K_a=\{x\in\K\ :\ ax=xa\}$. Montrer que $\K_a$ est un corps gauche dont $Z$ est un sous-corps. On note $d(a)$ la dimension de $\K_a$ en tant que $Z$-espace vectoriel. En étudiant la dimension de $\K$ en tant que $\K_a$ espace vectoriel, montrer que $d(a)$ divise $n$.
		\item En considérant les classes d'équivalence de la relation d'équivalence $x\sim y\iff\exists u\in\K\st,\ y=u\inv xu$ sur $\K\st$, montrer qu'il existe $d_1,\ldots,d_k$ des diviseurs stricts de $n$ tels que :
		$$q^n-1=q-1+\sum_{i=1}^{k}\frac{q^n-1}{q^{d_i}-1}.$$
		\item En déduire que $\Phi_n(q)$ divise $q-1$.
		\item En étudiant le module de $q-u$, où $u$ est une racine primitive $n$-ième de l'unité, puis le module de $\Phi_n(q)$, montrer que $n=1$ et conclure.
	\end{enumerate}
\end{pro}
\begin{sol}
	\begin{enumerate}
		\item Pour la première égalité, partitionner les racines $n$-ièmes de l'unité en les différents groupes de racines $d$-ièmes de l'unité, pour $d$ diviseur de $n$. Ensuite, récurrence sur $n$, en montrant que le théorème de division euclidienne s'adapte ici au polynômes à coefficients dans un anneau lorsque le coefficient dominant du diviseur est inversible (ici c'est $1$).
		\item Écrire $n=dk$ et $X^n-1=(X^d)^k-1^k=(X^d-1)\times(\cdots)$. Ensuite, dans le produit obtenu en $1$, sortir $\Phi_n$ et les $\Phi_k$ pour $k\mid d$.
		\item Prendre pour loi externe la restriction de la multiplication interne. Lorsque $\mathbb{M}$ est fini, $\mathbb{M}$ est de dimension finie sur $\L$, donc $\mathbb{M}$ est isomorphe à $\L^{\dim_\L(\mathbb{M})}$, si bien que :
		$$\card(\mathbb{M})=\card(\L)^{\dim_\L(\mathbb{M})}.$$
		\item Les vérifications sont immédiates. D'après la question précédente, $\K$ est un $Z$-espace vectoriel de dimension finie et $\card(\K)=q^n$.
		\item $\K_a$ est un sous-corps gauche de $\K$. On vérifie aisément que $Z$ est un sous-corps de $\K_a$. Par $3$, on a $\card(\K_a)=q^d$, puis $\card(\K)=(q^{d(a)})^k$ où $k$ est la dimension de $\K$ en tant que $\K_a$-espace vectoriel, et comme $q\geqslant 2$, on a $n=d(a)k$ donc $d(a)\mid n$.
		\item Si $x\in Z$, $\overline{x}=\{x\}$ ne possède qu'un élément. Si $x\notin Z$, il existe $y$ tel que $yx\ne xy$ donc $y\ne 0$ et $y\inv xy\ne x$ donc la classe de $x$ possède au moins deux éléments. Notons $\K_x\st=\K_x\setminus\{0\}$. $\K_x\st$ est un sous-groupe de $\K\st$ donc $q^{d(x)}-1$ divise $q^n-1$. On montre que le groupe quotient $\K\st /\K_x\st$ est équipotent à $\overline{x}$ via $u\K_x\st\mapsto u\inv x u$. Et comme $\overline{x}$ possède au moins deux éléments, $d(x)$ est un diviseur strict de $n$. \\
		On écrit alors $\K\st$ comme la partition des classes d'équivalences et on regarde les cardinaux pour obtenir l'égalité souhaitée. Peut-être mettre cette question en dernier ou la faire chercher en dernier s'il reste du temps.
		\item Par $1$, $\Phi_n$ divise $X^n-1$ et par $2$, $\Phi_n$ divise les $\displaystyle\frac{X^n-1}{X^{d_i}-1}$ (dans $\Z[X]$ à chaque fois), donc l'égalité de la question précédente permet de conclure.
		\item Si $n>1$, on voit aisément sur un dessin que pour toute racine primitive $n$-ième de l'unité $u$, on a $\lvert q-u\rvert>q-1$ (en effet, on rappelle que $q\geqslant 2$), donc si $n>1$, on a $\lvert\Phi_n(q)\rvert>(q-1)^\varphi(n)\geqslant q-1$. Or, $\lvert\Phi_n(q)\rvert\leqslant q-1$ d'après la question précédente, donc $n=1$, si bien que $\K=Z$ est commutatif !
	\end{enumerate}
\end{sol}

\begin{pro}[Autour de $\mathcal{SL}_n(\Z)$]
	Soit $G$ un groupe admettant une partie génératrice finie, et $H$ un groupe fini. On se donne $f$ un endomorphisme surjectif de $G$, et $g$ un morphisme de groupes de $G$ dans $H$.
	\begin{enumerate}
		\item Montrer qu'il y a un nombre fini de morphismes de $G$ dans $H$.
		\item Soit $a\in\ker(f)$. Montrer que pour tout $n\in\N$, il existe $b_n\in G$ tel que $a=f^n(b_n)$ et en déduire que $\ker(f)\subset\ker(g)$.
		\item On admet que $\mathcal{SL}_n(\Z)$ est un groupe admettant une partie génératrice finie. Montrer que tout endomorphisme de $\mathcal{SL}_n(\Z)$ est bijectif.
		\item On cherche à démontrer le résultat admis. Montrer que $\mathcal{SL}_n(\Z)$ est un groupe contenant les $T_{i,j}=I_n+E_{i,j}$ pour $i\ne j$. Calculer $T_{i,j}^k$ pour $k\in\Z$ et en déduire que $\mathcal{SL}_n(\Z)$ admet une partie génératrice finie.
	\end{enumerate}
\end{pro}
\begin{sol}
	\begin{enumerate}
		\item Un tel morphisme est caractérisé par les images des éléments de la partie génératrice. Or, la partie génératrice est finie, et pour chaque élément de cet ensemble, l'image est à choisir dans $H$ qui est fini, donc on a un nombre fini de possibilités.
		\item Si $n\in\N$, $f^n$ est surjectif, ce qui permet d'obtenir $b_n$.\\
		Ensuite, d'après la première question, il existe $n<m$ tel que $g\circ f^n=g\circ f^m$. Alors $g(a)=g\circ f^n(b_n)=g\circ f^m(b_n)=g\circ f^{m-n}(a)=e_G$.
		\item Soit $\psi$ un endomorphisme de $\mathcal{SL}_(\Z)$. Soit $p$ un nombre premier. On considère $\varphi_p:\mathcal{SL}_n(\Z)\to\mathcal{SL}_n(\Z/p\Z)$ la projection canonique. C'est aisément un morphisme de groupes et si $M\in\ker(\varphi_p)$, alors $M-I_n$ a tous ses coefficients divisibles par $p$. Ainsi, $\ker(\psi)\subset\ker(\varphi_p)$, et ceci pour tout $p$ premier. Par conséquent, si $M\in\ker(\psi)$, alors pour tout $p$ premier, tous les coefficients de $M-I_n$ sont divisibles par $p$, donc $M=I_n$ et $\psi$ est injectif.
		\item Pour la structure de groupe, utiliser la comatrice pour le passage à l'inverse. Le reste est aisé. Le fait que les $T_{i,j}$ y appartiennent est trivial. \\
		$T_{i,j}^k=I_n+kE_{i,j}$ (un binôme de Newton peut faire l'affaire pour s'éviter la récurrence). \\
		En utilisant l'algorithme d'Euclide colonne par colonne en multipliant à droite par des $T_{i,j}^k$, on ramène une matrice de $\mathcal{SL}_n(\Z)$ à une matrice dont la première ligne nulle, sauf le premier coefficient qui vaut le PGCD des coefficients initialement présents sur la ligne. \\
		En développant par rapport aux lignes, on observe que le PGCD des éléments d'une ligne vaut en fait $1$. \\
		Ensuite, on recommence avec le bloc $(n-1)\times (n-1)$ en bas à droite de la matrice provisoirement obtenue qui est dans $\mathcal{SL}_{n-1}(\Z)$. On se ramène par récurrence à une matrice triangulaire inférieure à coefficients entiers dont les termes diagonaux valent $1$.\\
		On annule ensuite les coefficients sous-diagonaux par opération (entières) sur les colonnes, et on s'est finalement ramené à $I_n$ en multipliant par des matrices $T_{i,j}^k$, si bien que les $T_{i,j}$ engendrent $\mathcal{SL}_n(\Z)$.
	\end{enumerate}
\end{sol}

\begin{pro}
	Soit $P=\displaystyle\sum_{k=0}^{n}a_kX^k\in\Z[X]$. On note $c(P)=\text{PGCD}(a_0,\ldots,a_n)$ le contenu de $P$.
	\begin{enumerate}
		\item Soit $(P,Q)\in\Z[X]^2$. Montrer que $c(PQ)=c(P)c(Q)$. Indication : on pourra se ramener au cas où $c(P)=c(Q)=1$.
		\item Montrer que l'on peut prolonger la fonction $c$ à $\Q[X]$ en conservant la propriété de la première question.
		\item Application 1 : retrouver le fait que les racines rationnelles d'un polynôme unitaire à coefficients entiers sont entières.
		\item Application 2 : soit $P\in\Z[X]$ non constant. Montrer que $P$ est irréductible dans $\Z[X]$ si, et seulement s'il est irréductible dans $\Q[X]$ et $c(P)=1$.
		\item Application 3 : montrer que le polynôme minimal d'une matrice de $\mathcal{M}_n(\Z)$, considérée comme une matrice de $\mathcal{M}_n(\Q)$, est à coefficients entiers.
	\end{enumerate}
\end{pro}
\begin{sol}
	\begin{enumerate}
		\item On procède en deux étapes.
		\begin{itemize}
			\item Supposons $c(A)=c(B)=1$. Par l'absurde, supposons $c(AB)\ne 1$. Prenons alors $p$ un nombre premier qui divise $c(AB)$. Alors, en passant au quotient dans $\Z/p\Z[X]$, on a $\overline{0}=\overline{AB}=\overline{A}\times\overline{B}$. Puisque $\Z/p\Z$ est un corps, $\Z/p\Z[X]$ est intègre donc $\overline{A}=\overline{0}$ ou $\overline{B}=\overline{0}$, et donc $p$ divise $c(A)$ ou $c(B)$, ce qui est absurde. Donc $c(AB)=1$.
			\item Dans le cas général, on écrit $A=c(A)A'$ et $B=c(B)B'$ avec $c(A')=c(B')=1$ (en effet, $c(nP)=\lvert n\rvert c(P)$). Alors, par le cas précédent : $$c(AB)=c\big(c(A)c(B)A'B'\big)=c(A)c(B)c(A'B')=c(A)c(B)$$
		\end{itemize}
		\item Poser le seul prolongement naturel (avec le PPCM ou en faisant les produits) et vérifier que cela est bien défini et conserve la propriété précédente.
		\begin{rmq}
			Il est important de noter que si $P\in\Q[X]$, alors $P/c(P)$ est un polynôme à coefficients entiers de contenu $1$. C'est l'argument essentiel à utiliser pour les trois applications suivantes.
		\end{rmq}
	\end{enumerate}
\end{sol}

\section{Exercices}

\begin{ex}
	Soit $p$ un nombre premier. Montrer qu'il existe $n\in\N$ tel que $2^n\equiv n\ [p]$.
\end{ex}
\begin{sol}
	Prendre $n=2$ si $p=2$. On suppose $p\geqslant 3$. Alors $2^{p-1}\equiv 1\ [p]$ d'après le petit théorème de Fermat. Donc si on écrit la division euclidienne d'un potentiel $n$ solution par $p-1$ sous la forme $n=q(p-1)+r$, on a $2^n\equiv 2^r\ [p]$, et donc $2^r\equiv -q+r\ [p]$. Il suffit de choisir $q$ et $r$ vérifiant ceci pour conclure. En imposant $r=0$, on peut par exemple prendre $q=p-1$, et vérifier que cela convient.
\end{sol}

\begin{ex}
	Soit $A$ un anneau commutatif fini. On note $N$ l'ensemble de ses éléments nilpotents. Montrer que $\card(N)$ divise $\card(A)$ et $\card(A^\times)$.
\end{ex}
\begin{sol}
	$(N,+)$ est un sous-groupe de $(A,+)$, puis utiliser le théorème de Lagrange. \\
	Pour la deuxième, l'ensemble des $1-n$ pour $n\in N$ (en bijection avec $N$) est un sous-groupe multiplicatif de $(A^\times,\times)$, puis on utilise à nouveau le théorème de Lagrange.
\end{sol}

\begin{ex}
	Soit $A$ un anneau tel que : $\forall (x,y)\in A^2,\ (xy=yx\ \text{ou}\ xy=-yx)$. Montrer que $A$ est commutatif.
\end{ex}
\begin{sol}
	Attention, le symbole $+$ ou $-$ dépend du couple $(x,y)$ !!! Soit $x\in A$. On pose :
	$$A_+(x)=\{y\in A\ :\ xy=yx\}\quad\text{et}\quad A_-(x)=\{y\in A\ :\ xy=-yx\}.$$
	On montre aisément que ces deux ensembles sont des sous-groupes additifs de $(A,+)$.\\
	Mais $A=A_+(x)\cup A_-(x)$ par hypothèse. Or, l'union de deux sous-groupes est un sous-groupe si, et seulement si, un des sous-groupes est inclus dans l'autre. Ici, on a donc $A_+(x)=A$ ou $A_-(x)=A$.\\
	Dans le premier cas, on a gagné.\\
	Dans le deuxième cas, on en déduit notamment que $1_A\in A_-(x)$ donc $x=-x$. Puis si $y\in A=A_-(x)$, alors $xy=-yx=y(-x)=yx$ donc on a aussi gagné !
\end{sol}

\begin{ex}[Module d'une somme de Gauss]
	Soit $p$ un nombre premier impair. On pose :
	$$S_p=\sum_{k=0}^{p-1}\exp\left(\frac{2i k^2\pi}{p}\right).$$
	\begin{enumerate}
		\item Montrer que :
		$$\begin{array}{rcl}
			(\Z/p\Z)^2&\longrightarrow&(\Z/p\Z)^2\\
			(x,y)&\longmapsto&(x+y,x-y)
		\end{array}$$ est bijective.
		\item Montrer que $\lvert S_p\rvert=\sqrt{p}$.
	\end{enumerate}
\end{ex}
\begin{sol}
	\begin{enumerate}
		\item C'est un morphisme de groupes additifs. On montre que le noyau est réduit à $\{(0,0)\}$, pour cela, on utilise le fait que $2$ est inversibles car $p$ est premier impair.
		\item Calculer $\lvert S_p\rvert^2=S_p\overline{S_p}$. On trouve :
		$$\sum_{0\leqslant k,l\leqslant p-1}\exp\left(\frac{2i(k-l)(k+l)\pi}{p}\right).$$
		On peut remplacer $k$ et $l$ par leurs classes modulo $p$, ce qui ne change rien, puis utiliser le changement de variable obtenu à la question précédente pour se ramener à des sommes géométriques. Seule une somme géométrique est non nulle et vaut $p$. On en déduit le résultat.
	\end{enumerate}
\end{sol}

%\begin{ex}
%	Soit $G$ un groupe de cardinal $2n$.
%	\begin{enumerate}
%		\item Soit $H$ un sous-groupe de $G$ de cardinal $n$. Montrer que l'application :
%		$$\begin{array}{lrcl}
%			\varphi:&G&\longrightarrow&\U_2\\
%			&g&\longmapsto&\displaystyle\begin{cases}
%				1&\text{si }g\in H\\
%				-1&\text{si }g\notin H
%			\end{cases}
%		\end{array}$$ définit un morphisme de groupes de $(G,\cdot)$ dans $(\U_2,\times)$.
%		\item Pour $d\in\N\st$, combien $(\Z/2\Z)^d$ possède-t-il de sous-groupes de cardinal $2^{d-1}$ ?
%		\item Montrer que $G$ ne peut pas posséder exactement deux sous-groupes de cardinal $n$.
%	\end{enumerate}
%\end{ex}
%\begin{sol}
%	\begin{itemize}
%		\item Faire une table de multiplication peut aider à clarifier le raisonnement.
%		\item On peut voir $(\Z/2\Z)^d$ comme un $\Z/2\Z$-espace vectoriel de dimension $d$. Un tel sous-groupe en est alors un hyperplan. Il y a autant d'hyperplans que de droites, donc il y en a ici $2^d-1$.
%		\item La question $1$ permet de montrer qu'à tout sous-groupe de $G$ de cardinal $n$ correspond un unique morphisme de groupe de $G$ dans $\Z/2\Z$. Or, l'espace de ces morphismes est un $\Z/2\Z$-espace vectoriel de dimension finie.
%	\end{itemize}
%\end{sol}


\chapter{}

\section{Problèmes}

\section{Exercices}

\end{document}